\scanpage{29}
\deleted{\unsure{SPS|SQS} de plus que $\pi_1(X)$ a un centre trivial. En car. $p>0$, notons qu'il suffit de travailler avec les groupes fondamentaux premiers à $p$ -- et $\unsure{\overline m}$ sur les localisations en $\ell\ne p$ des groupes fondamentaux.}
\mspage{14}

La situation $X,Y,f,g$ provient, par extension du corps de base, d'une situation analogue sur un corps $K$ \inserted{extension} de type fini de $\mathbb Q$. \deleted{Choisissons une} \inserted{Soit $\overline K$ la} clôture algébrique \deleted{$\unsure{\overline K|\overline Q}$} de $K$ dans $k$, \deleted{\unsure{des $k$ l'}} \inserted{et $\overline X,\overline Y,\overline f,\overline g$ se} déduite de \inserted{$X,Y$,} $f,g$ par extension du corps de base de $\unsure{K|Q}$ à $\overline K$. On a donc $\pi_1(\overline f)$, $\pi_1(\overline g)$ satisfaisant la condition d) avec de plus $\pi_1(\overline X)$ à centre trivial.

Mais ces hypothèses impliquent que les extensions $E_{\inserted{X/K}}\deleted{(X)}\inserted{=\pi_1(X)}$, $E_{\inserted{Y/K}}\deleted{(Y)}\inserted{=\pi_1(Y)}$ de $E_{\overline K/K}$ par $\pi_1(\overline X)$, resp. $\pi_1(\overline Y)$, ainsi que les hom. $\pi_1(f)$, $\pi_1(g)$ induits, sont reconstruits à partir de \inserted{$\pi_1(\overline X)$, $\pi_1(\overline Y)$, de} $\pi_1(\overline f)$, $\pi_1(\overline g)$ et de l'action extérieure de $E_{\overline K/K}$ sur ces groupes. On va montrer maintenant le
\margin{28}
